This chapter presents the experimental results and evaluates the behavior of the decentralized learning system under a targeted backdoor attack. The analysis focuses on how different network architectures affect both the learning of legitimate data and the propagation of the malicious infection. The evaluation is divided into several phases: establishing the baseline resilience of standard topologies, analyzing the attacker's trade-off between effectiveness and stealthiness, evaluating community-based structures (SBM), and testing the system under scaled and more realistic conditions.


\section{Baseline Topologies Resilience}
The first phase of the evaluation analyzed the propagation of the continuous backdoor attack across three fundamental communication structures: Fully Connected, Ring, and Star. To accurately measure the long-term assimilation of the payload, the network consisted of 5 nodes operating over 100 communication rounds. The adversary (Client 1) employed the optimized stealth configuration defined in Section 4: a 20\% local poisoning ratio and the topological boosting heuristic of $\alpha = d + 1$. 

The most immediate and significant observation across all three baseline topologies is the complete preservation of the attacker's stealth. As shown in the top subplots of Figures \ref{fig:metrics_fullyconnected}, \ref{fig:metrics_ring}, \ref{fig:metrics_star_peripheral} and \ref{fig:metrics_star_hub}, the Clean Accuracy of the malicious node remains consistently above 0.95, mirroring the performance of the honest majority. By abandoning the brute-force single-shot strategy and utilizing the continuous, degree-scaled heuristic, the attacker successfully avoids Catastrophic Forgetting, rendering the malicious node virtually indistinguishable to standard anomaly detection systems.

\begin{figure}[t]
    \centering
    \subfloat[Fully Connected\label{fig:metrics_fullyconnected}]{%
        \includegraphics[width=0.45\textwidth]{"../results/outputs/2026-05-04_13-19-fc_5/metrics_per_node_FullyConnected.png"}%
    }\hfill
    \subfloat[Ring Topology\label{fig:metrics_ring}]{%
        \includegraphics[width=0.45\textwidth]{"../results/outputs/2026-05-04_15-41-ring_5/metrics_per_node_ring.png"}%
    }
    \caption{Clean Accuracy and Attack Success Rate (ASR) over 100 rounds for the Fully Connected and Ring topologies.}
    \label{fig:metrics_baseline_fc_ring}
\end{figure}


\begin{figure}[t]
    \centering
    \subfloat[Star Topology — Peripheral Attacker\label{fig:metrics_star_peripheral}]{%
        \includegraphics[width=0.45\textwidth]{"../results/outputs/2026-05-04_16-25-star_5/metrics_per_node_star.png"}%
    }\hfill
    \subfloat[Star Topology — Hub Attacker\label{fig:metrics_star_hub}]{%
        \includegraphics[width=0.45\textwidth]{"../results/outputs/2026-05-12 - 17_26 - star_5 - client0/metrics_per_node.png"}%
    }
    \caption{Clean Accuracy and ASR over 100 rounds for the Star topology under two distinct attacker placements.}
    \label{fig:metrics_star_combined}
\end{figure}

However, while stealth was universally maintained, the Attack Success Rate (ASR) varied drastically depending on the structural placement of the nodes:

\textbf{Fully Connected Topology (Figure \ref{fig:metrics_fullyconnected}):}
In the maximally dense Fully Connected network, the lack of communication barriers allowed the malicious updates to slowly but relentlessly infiltrate the global consensus. While the honest nodes successfully resisted the backdoor for the first 30 rounds, the continuous injection strategy eventually overwhelmed them. Starting around round 45, the ASR for all honest nodes (Clients 0, 2, 3, and 4) began a unified, steady climb, eventually reaching approximately 0.80 by round 100. This demonstrates that in highly connected, symmetric networks, the $\alpha = d+1$ heuristic perfectly overcomes neighborhood dilution, resulting in a system-wide infection.

\textbf{Ring Topology (Figure \ref{fig:metrics_ring}):}
Conversely, the Ring topology demonstrated a strong structural resistance to rapid propagation. Due to the strict neighborhood constraints ($d=2$), the backdoor was forced to travel sequentially. The ASR graph reveals a clear staggered infection: the attacker's immediate neighbors (Clients 0 and 2) slowly climbed to an ASR of approximately 0.40 by the end of the simulation. However, the distant nodes (Clients 3 and 4) were largely shielded by the topological distance, maintaining an ASR below 0.20 for over 90 rounds. The structural bottleneck of the Ring effectively trapped the infection, proving that sparsity natively buys the network time against continuous poisoning.

\textbf{Star Topology (The Center of Power):}
To fully understand the vulnerabilities of localized centralization within a peer-to-peer framework, the Star topology was evaluated under two distinct attacker placements: an attack from an isolated edge (Peripheral Attacker) and an attack from the structural center (Hub Attacker).

When the adversary is placed at the periphery (Figure \ref{fig:metrics_star_peripheral}), the topology exhibits a formidable defense mechanism. The attacker (Client 1) attempts to push its backdoor, but the central hub (Client 0) continuously averages the poisoned weights with clean updates from the other three isolated nodes. This central aggregation acts as an impenetrable "consensus wall," completely neutralizing the malicious parameters. The ASR for the hub and all honest peripheral nodes remains flat near 0.0 for the entire simulation.

However, when the adversary compromises the central Hub itself (Figure \ref{fig:metrics_star_hub}), the network becomes exceptionally fragile. Because Client 0 connects to all nodes—and the peripheral nodes have no lateral connections to one another—the attacker possesses a direct, unfiltered broadcast channel. As demonstrated by the ASR graph, the attacker quickly embeds the backdoor into its local model and systematically transmits it. Without any alternative honest neighbors to dilute the payload, the peripheral nodes are defenseless. By round 60, the ASR for all clients surges past 0.80, resulting in a devastating network-wide compromise while the attacker perfectly preserves its stealth (Clean Accuracy > 0.95).

These baseline results definitively prove that an attack's effectiveness is not strictly tied to the payload itself, but is governed entirely by the architectural constraints of the decentralized graph.

\section{Solving the Effectiveness vs. Stealthiness Trade-off}
A crucial observation across all evaluated topologies is the definitive resolution of the effectiveness versus stealthiness trade-off. Traditionally, high-intensity model poisoning—characterized by aggressive boosting factors and high local poisoning ratios—risks triggering Catastrophic Forgetting \cite{french1999catastrophic}. In such scenarios, the excessive parameter distortion causes the local model's Clean Accuracy to collapse, rendering the malicious node easily detectable by standard anomaly detection systems.

However, the implementation of our proposed Continuous Poisoning strategy (utilizing a 20\% poison ratio) combined with the topologically-derived boosting heuristic ($\alpha = d + 1$) successfully eliminates this vulnerability. As demonstrated in the top subplots of all figures in this chapter, the compromised node invariably maintained a Clean Accuracy above 0.90, mirroring the performance of honest nodes. By carefully pacing the injection across 100 to 160 rounds and scaling the payload mathematically to the network's local dilution rate, the adversary maximized its effectiveness while remaining completely stealthy across every tested architectural constraint.

\section{The Topological Firewall (SBM Peripheral Attacker)}
To evaluate the resilience of realistic, community-based networks, we simulated an 8-node Stochastic Block Model (SBM) divided into two distinct communities (Community A: Clients 0-3; Community B: Clients 4-7) connected by a single gateway link between Client 3 and Client 4. The simulation spanned 160 rounds to test long-term boundary resilience.

\begin{figure}[t]
	\centering
	\subfloat[SBM Peripheral Attacker\label{fig:metrics_sbm_peripheral}]{%
		\includegraphics[width=0.45\textwidth]{"../results/outputs/2026-05-05_15-37-sbm_8/metrics_per_node.png"}%
	}\hfill
	\subfloat[SBM Gateway Attacker\label{fig:metrics_sbm_gateway}]{%
		\includegraphics[width=0.45\textwidth]{"../results/outputs/2026-05-05_16-52-sbm_8-client3/metrics_per_node_sbm.png"}%
	}
	\caption{Clean Accuracy and ASR over 160 rounds for the Stochastic Block Model under two distinct attacker placements.}
	\label{fig:metrics_sbm_combined}
\end{figure}

In the first scenario (Figure \ref{fig:metrics_sbm_peripheral}), the attacker is placed at Client 1, deep within Community A. The results demonstrate a clear "echo chamber" effect. Because Client 1 is well-connected locally but isolated from the broader network, the backdoor successfully circulates within Community A, pushing the ASR of its immediate neighbors (Clients 0 and 2) to approximately 0.20. 

However, the attack entirely fails to cross the community boundary. The ASR for all nodes in Community B (Clients 4, 5, 6, and 7) remains flat at 0.0 for the entire 160 rounds. This proves that a sparse gateway link natively functions as a physical, topological firewall. The localized consensus of Community B is too mathematically dense for a peripheral infection to penetrate via a single bridging connection.

\section{The Bottleneck Breach Attempt (SBM Gateway Attacker)}
To push the network to its absolute limit, the final experiment shifted the compromised node to the structural bottleneck itself. In this scenario (Figure \ref{fig:metrics_sbm_gateway}), Client 3 (the gateway node of Community A) acts as the adversary, utilizing a higher dynamic boosting factor ($\alpha = d+1$) commensurate with its higher connectivity.

The change in attacker placement yielded two profound discoveries:
\begin{itemize}
    \item \textbf{Amplified Local Infection:} Because the attacker was positioned at a highly central point within Community A, the local infection rate skyrocketed. Clients 0, 1, and 2 experienced an ASR surge to roughly 0.50—more than double the effectiveness of the peripheral attack. Controlling the gateway allowed the adversary to dominate its own community.
    \item \textbf{The "Consensus Sink" Defense:} Remarkably, despite the attacker having a direct edge to Community B, the broader network still survived. The ASR for Clients 4 through 7 remained near 0.0. 
\end{itemize}

This ultimate failure to bridge the gap highlights the profound defensive power of Decentralized Federated Learning. When the gateway node of Community B (Client 4) received the highly poisoned parameters from Client 3, it immediately aggregated them with the clean updates from its internal neighbors (Clients 5, 6, and 7). Community B effectively acted as a massive "consensus sink," structurally washing out the backdoor before it could take root. 

Ultimately, these experiments prove that while continuous, topology-aware backdoor attacks can successfully evade detection and dominate dense local clusters, decentralized community boundaries provide a natural, highly robust defense against global network corruption.
